\chapter{Vitamin B3 (Niacin) Deficiency --- Pellagra}
\index{Vitamin B3 (Niacin) Deficiency --- Pellagra}
\label{sec:Vitamin B3 Deficiency}

\section{Overview}


\begin{itemize}[noitemsep]
  \item Niacin (vitamin B3) deficiency (pellagra) is a vitamin deficiency disorder classically associated with diets based on maize (corn), which contains niacin in a bound, non-bioavailable form unless treated with alkali (nixtamalization)
  \item It also occurs in alcoholism, carcinoid syndrome, and with isoniazid therapy.
\end{itemize}
\section{Causes}
This condition happens because of deficiency of niacin, a precursor for nicotinamide adenine dinucleotide (NAD) and NADP, essential for redox reactions and cellular energy metabolism. Metabolic disorders arise from dysregulation of normal biochemical pathways. This may result from enzyme deficiencies, hormonal imbalances, or lifestyle factors that overwhelm normal homeostatic mechanisms.   
\section{Risk Factors}
\begin{itemize}[noitemsep]
  \item Genetic predisposition and family history
  \item Advancing age
  \item Lifestyle factors (diet, exercise, smoking, alcohol)
  \item Environmental exposures
  \item Underlying medical conditions
  \item Medication use
\end{itemize}
\section{Signs and Symptoms}
\subsection{Common symptoms}
\begin{itemize}[noitemsep]
  \item You may experience the classic triad of dermatitis, dementia, and diarrhea (the ``three D's'')
  \item Dermatitis is photosensitive, symmetric, with well-demarcated, erythematous, hyperpigmented, scaling plaques in sun-exposed areas (Casal necklace)
  \item Digestive tract involvement includes diarrhea, glossitis, and stomatitis
  \item Neuropsychiatric manifestations include depression, apathy, confusion, memory loss, and psychosis.
  \item Pain or discomfort in the affected area
  \end{itemize}
\subsection{Emergency warning signs}
\begin{itemize}[noitemsep]
  \item Severe or worsening symptoms of vitamin b3 (niacin) deficiency --- pellagra require immediate medical evaluation
\end{itemize}
\section{Progression and Stages}
The condition may progress through different stages of severity over time. Early stages often involve mild or subtle symptoms that may be overlooked. Moderate stages involve more noticeable symptoms affecting daily function. Advanced stages may involve significant impairment and complications.
\section{Complications}
\begin{itemize}[noitemsep]
  \item Disease progression leading to worsening symptoms
  \item Secondary infections or injuries
  \item Side effects of treatment
  \item Reduced quality of life
  \item Emotional and psychological impact
\end{itemize}
\section{Diagnosis}
Doctors usually diagnose this based on your symptoms and a physical exam.

\begin{itemize}[noitemsep]
  \item Comprehensive medical history and physical examination
  \item Laboratory tests to assess organ function
  \item Imaging studies to visualize affected structures
  \item Specialized testing depending on the affected system
  \item Consultation with relevant specialists
\end{itemize}
\section{Prevention}
\begin{itemize}[noitemsep]
  \item Maintain a healthy lifestyle with balanced diet and exercise
  \item Avoid known risk factors
  \item Attend regular medical check-ups for early detection
  \item Follow recommended screening guidelines
\end{itemize}
\section{Treatment Options}
Treatment focuses on nicotinamide (preferred over niacin to avoid flushing) 100--500 mg orally daily. Dietary modifications and nutritional counseling Pharmacotherapy to correct metabolic abnormalities Hormone replacement therapy when indicated Regular monitoring of metabolic parameters Weight management programs Exercise prescription Management of associated conditions Patient education for self-management

\begin{itemize}[noitemsep]
  \item Medications to manage symptoms and address underlying mechanisms
  \item Lifestyle modifications and self-care measures
  \item Physical therapy and rehabilitation
  \item Surgical intervention when indicated
  \item Regular monitoring and follow-up care
\end{itemize}
\section{Living with It}
\begin{itemize}[noitemsep]
  \item Follow your treatment plan as prescribed
  \item Maintain regular follow-up with your healthcare provider
  \item Make necessary lifestyle adjustments
  \item Seek support from family, friends, or support groups
  \item Stay informed about your condition
\end{itemize}
\section{Outlook}
Most people recover well with treatment

\begin{itemize}[noitemsep]
  \item Dermatitis and digestive tract symptoms improve within days
  \item Nervous system symptoms resolve more slowly.
\end{itemize}
